\chapter{Descripción general}
\label{ch:descripcion}

\section{El problema de Job Shop Scheduling}

Formalmente, una instancia del \gls{jsp} está compuesta por un conjunto de \(n\) trabajos y un conjunto de \(m\) máquinas. Cada trabajo \(J_i\), con \(i=1,\ldots,n\), debe ejecutar una secuencia ordenada de operaciones
\[
    O_{i1} \rightarrow O_{i2} \rightarrow \cdots \rightarrow O_{im_i},
\]
donde \(m_i\) es el número de operaciones del trabajo \(i\). Cada operación \(O_{ij}\) requiere la máquina \(m_{ij}\) durante un tiempo de procesamiento \(p_{ij}\) y no puede interrumpirse una vez iniciada. La secuencia de operaciones de cada trabajo representa su ruta tecnológica y debe respetarse estrictamente.

Sea \(s_{ij}\) el instante de inicio de la operación \(O_{ij}\). Las restricciones de precedencia dentro de cada trabajo se expresan como
\[
    s_{i,j+1} \ge s_{ij} + p_{ij}, \qquad j=1,\ldots,m_i-1.
\]
Además, dos operaciones que utilizan la misma máquina no pueden superponerse en el tiempo. Si \(O_{ij}\) y \(O_{kl}\) requieren la misma máquina, entonces debe cumplirse una de las dos condiciones siguientes:
\[
    s_{ij} + p_{ij} \le s_{kl}
    \qquad \text{o} \qquad
    s_{kl} + p_{kl} \le s_{ij}.
\]
El instante de término del trabajo \(i\) es \(C_i = s_{i,m_i} + p_{i,m_i}\), y el makespan se define como
\[
    C_{\max} = \max_{i=1,\ldots,n} C_i.
\]
El objetivo del problema es encontrar una programación factible que minimice \(C_{\max}\). En este informe se utilizan tres instancias de complejidad creciente: \(6\times6\), \(15\times10\) y \(50\times10\), donde el primer valor corresponde al número de trabajos y el segundo al número de máquinas.

\section{Metaheurística seleccionada}

\subsection{Tipo de metaheurística}

El recocido simulado es una metaheurística de trayectoria, también denominada de solución única. A diferencia de las metaheurísticas poblacionales, como los algoritmos genéticos o la optimización por enjambre de partículas, el recocido simulado mantiene una sola solución actual y la modifica iterativamente mediante un operador de vecindario \parencite{talbi2009}. Su funcionamiento se inspira en el proceso físico de recocido de los metales: a altas temperaturas los átomos pueden reorganizarse y aceptar configuraciones de mayor energía; a medida que la temperatura disminuye, el sistema converge hacia configuraciones de menor energía \parencite{kirkpatrick1983}.

En el contexto de optimización, la ``energía'' corresponde al valor de la función objetivo y la ``temperatura'' controla la probabilidad de aceptar soluciones peores. Si \(x\) es la solución actual, \(x'\) es un vecino y \(\Delta = f(x') - f(x)\), el criterio de aceptación de Metropolis es
\[
    P(\text{aceptar } x') =
    \begin{cases}
        1, & \text{si } \Delta < 0, \\[0.2cm]
        e^{-\Delta/T}, & \text{si } \Delta \ge 0,
    \end{cases}
\]
donde \(T>0\) es la temperatura actual. Cuando \(T\) es alta, la probabilidad de aceptar un empeoramiento es considerable y el algoritmo explora ampliamente el espacio de búsqueda. Cuando \(T\) se aproxima a cero, la probabilidad de aceptar soluciones peores disminuye y el algoritmo se comporta como una búsqueda local descendente.

\subsection{Ventajas y desventajas}

Las principales ventajas del recocido simulado son:

\begin{itemize}
    \item Es sencillo de implementar y requiere pocos parámetros: temperatura inicial, temperatura mínima, factor de enfriamiento y criterio de parada.
    \item Permite escapar de óptimos locales gracias a la aceptación probabilística de soluciones peores.
    \item Es un algoritmo \emph{anytime}: en cualquier momento de la ejecución se dispone de una solución factible y de la mejor solución encontrada hasta ese instante.
    \item Se adapta con facilidad a problemas combinatorios en los que es posible definir un vecindario y una función objetivo.
    \item Tiene un bajo costo computacional por iteración cuando la evaluación del vecindario puede realizarse de forma eficiente.
\end{itemize}

Sus principales desventajas son:

\begin{itemize}
    \item No garantiza encontrar la solución óptima ni permite certificar optimalidad.
    \item Su desempeño es sensible a la calibración de la temperatura inicial, la temperatura mínima y el factor de enfriamiento.
    \item Un enfriamiento demasiado rápido puede conducir a un óptimo local de baja calidad; un enfriamiento demasiado lento puede consumir un tiempo excesivo.
    \item Mantiene una sola solución, por lo que no aprovecha la diversidad de una población.
    \item La calidad de la solución depende del vecindario elegido y de la representación utilizada.
\end{itemize}

\subsection{Razón de la elección}

Se eligió el recocido simulado por tres razones principales. Primero, el problema estudiado es de naturaleza combinatoria y el recocido simulado ha demostrado ser eficaz en problemas de programación de la producción \parencite{pinedo2016}. Segundo, el algoritmo es fácil de explicar, implementar y reproducir, lo que facilita la entrega y la presentación del trabajo. Tercero, el grupo ya contaba con una implementación base en Rust, lo que permitió concentrar el esfuerzo en la representación de soluciones, la decodificación y la experimentación, en lugar de partir desde cero.

Además, el recocido simulado ofrece un buen equilibrio entre calidad y costo computacional para las tres instancias propuestas. Como se muestra en el Capítulo~\ref{ch:resultados}, con un presupuesto de aproximadamente 27.625 iteraciones por ejecución se alcanzó el óptimo en las tres instancias, lo que confirma que la elección fue adecuada para el alcance del trabajo.
